% Generated draft: requires mathematical and source review.
\documentclass[11pt]{article}
\usepackage[margin=25mm]{geometry}
\usepackage{fontspec}
\setmainfont{Latin Modern Roman}
\usepackage{amsmath,amssymb,mathtools,unicode-math}
\setmathfont{Latin Modern Math}
\usepackage{enumitem,needspace,xurl,hyperref}
\setlength{\parindent}{0pt}
\setlength{\parskip}{5pt}
\setlength{\emergencystretch}{3em}
\newcommand{\N}{\mathbb{N}}
\newcommand{\Z}{\mathbb{Z}}
\newcommand{\Q}{\mathbb{Q}}
\newcommand{\R}{\mathbb{R}}
\newcommand{\C}{\mathbb{C}}
\newcommand{\F}{\mathbb{F}}
\newcommand{\A}{\mathbb{A}}
\newcommand{\PP}{\mathbb{P}}
\newcommand{\E}{\mathbb{E}}
\newcommand{\eps}{\varepsilon}
\newcommand{\dd}{\,\mathrm{d}}
\begin{document}
\section*{UnsolvedMath ID: 2306099}
\textbf{Source problem code:} catalogue entry
\subsection{Definitions and mathematical statement}
Consider normalized holomorphic univalent functions \(f(z)=z+\sum_{n\ge2}a_nz^n\) on \(\mathbb D\), with convex image. The source defines \(CV(R_1,R_2)\) by requiring a boundary quantity called curvature \(\rho\) to satisfy
\[
R_1\le\rho\le R_2\quad\text{on }f(\{|z|=1\}).
\]
Its concrete extremal problem is to determine the sharp coefficient bound
\[
B(R_1,R_2)=\sup\{|a_2|:f\in CV(R_1,R_2)\}
=\tfrac12\sup\{|f''(0)|:f\in CV(R_1,R_2)\},
\]
for the intended parameter range in which the class is nonempty. Sharpness means the least valid upper bound; attainment is a separate issue.

The boundary and curvature conventions are missing. For a regular \(C^2\) boundary parametrization \(h(t)=f(e^{it})\), geometric curvature is
\[
\kappa(t)=\frac{|\operatorname{Im}(\overline{h'(t)}h''(t))|}{|h'(t)|^3}.
\]
If \(\rho\) means radius of curvature instead, then \(\rho=1/\kappa\), giving different parameter inequalities and a different extremal class. The source explicitly says curvature, but its notation and intended convention must be checked. Neither a smooth boundary extension nor a generalized curvature notion is supplied, so the expression \(f(\{|z|=1\})\) is not automatically defined for an arbitrary disk map. The values and admissible range of \(R_1,R_2\) are also unstated.

Goodman's source request includes broader study of this class, with the sharp \(|a_2|\) bound as the explicit unresolved target. The supplied OPEN-TRIAGE summary reports open status in the 2018 edition and no definitive resolution located; no new bound or source correction is asserted.
\subsection{Short English statement}
Find the best possible second Taylor-coefficient bound for normalized convex disk maps with two-sided boundary-curvature constraints. The curvature-versus-radius convention, boundary regularity, and admissible parameters require review.
\subsection{Sources}
\begin{itemize}
\item Local UnsolvedMath registry, permanent ID \(2306099\), exact source code \texttt{catalogue entry}, ``Research Problems in Function Theory --- Problem 6.99.''
\item Hayman, \emph{Research Problems in Function Theory} (2018), Problem 6.99, attributed to A. W. Goodman; \url{https://arxiv.org/abs/1809.07200}. Source-TeX extraction; recorded access 2026-07-29.
\item The statement says ``see Goodman and'' with incomplete/missing references. The supplied OPEN-TRIAGE summary gives no separate current-resolution citation.
\end{itemize}
\end{document}
